% GENERATED FILE. Edit canonical lessons, then run npm run generate:books.
% canonical-source-hash: fnv1a64:e4933cec624b98f0
% canonical-lessons: MW-C273-chup, MW-C273-isaro, MW-C273-bhed, MW-C273-chhipavno, MW-C273-vachan

\chapter{Hush, a Secret}
\label{ch:mw-hush-a-secret}
\hlchaptermodality{273}

\begin{chapteropening}
\textbf{By the end of this chapter:} \emph{I can ask for quiet and keep a secret.}

Complete the last lesson of chapter 273: I can ask for quiet and keep a secret.
\end{chapteropening}

\section[\texorpdfstring{\mw{चुप} (chup)}{चुप (chup)}]{\mw{चुप} (chup) — silent}

\label{lesson:MW-C273-chup}



\begin{quote}
Before the new one: say the Marwadi for of course, then the Marwadi for agreed.
\end{quote}

\subsection*{You'll want to know: \mw{चुप}}
\textbf{\mw{चुप}} — \emph{chup} — \textquotedblleft{}silent\textquotedblright{}.

\begin{quote}
\mw{चुप} \mw{रैवो।}
\end{quote}

\begin{quote}
\emph{chup raivo.}
\end{quote}

\begin{quote}
Keep quiet.
\end{quote}

\begin{cousinweb}
\textbf{\mw{चुप}} (\emph{chup}) is inherited from a reconstructed \emph{*cuppa}, \textquotedblleft{}silent\textquotedblright{}; CDIAL gives no Marwari form in this entry, but lists Hindi \emph{cup}, \textquotedblleft{}silent\textquotedblright{} (CDIAL 4864).
\end{cousinweb}

\begin{grammarlens}[title={What you've built}]
One more word for asking and answering.
\end{grammarlens}

\subsection*{Your turn}
\begin{itemize}
\raggedright
  \item \emph{Say it:} \emph{chup}
  \item \emph{Say it:} \emph{chup}, once more
  \item \emph{Say it:} \emph{chup raivo.} — Keep quiet.
  \item \emph{Recall:} say the Marwadi for of course, then the Marwadi for agreed, then say \emph{chup} again
\end{itemize}

\begin{tcolorbox}[breakable,colback=teal!4,colframe=teal!35!black,title={Before you move on}]
What does \mw{चुप} mean? (\textquotedblleft{}Silent\textquotedblright{}.) Say it once more.
\end{tcolorbox}
\section[\texorpdfstring{\mw{इसारो} (isāro)}{इसारो (isāro)}]{\mw{इसारो} (isāro) — a gesture}

\label{lesson:MW-C273-isaro}



\begin{quote}
Before the new one: say the Marwadi for agreed, then the Marwadi for silent.
\end{quote}

\subsection*{You'll want to know: \mw{इसारो}}
\textbf{\mw{इसारो}} — \emph{isāro} — \textquotedblleft{}a gesture\textquotedblright{}.

\begin{quote}
\mw{मास्टर} \mw{इसारो} \mw{करै} \mw{है} \mw{अर} \mw{छोरा} \mw{चुप} \mw{है।}
\end{quote}

\begin{quote}
\emph{māsṭar isāro karai hai ar chhorā chup hai.}
\end{quote}

\begin{quote}
The teacher makes a sign and the boys are quiet.
\end{quote}

A masculine noun.

\begin{cousinweb}
\textbf{\mw{इसारो}} (\emph{isāro}) is usually traced, through Persian \emph{ishāra}, to Arabic \emph{ishāra}, \textquotedblleft{}a sign\textquotedblright{}; Marwadi gives it the -\mw{ो} ending of a masculine noun.
\end{cousinweb}

\begin{grammarlens}[title={What you've built}]
One more word for asking and answering.
\end{grammarlens}

\subsection*{Your turn}
\begin{itemize}
\raggedright
  \item \emph{Say it:} \emph{isāro}
  \item \emph{Say it:} \emph{isāro}, once more
  \item \emph{Say it:} \emph{māsṭar isāro karai hai ar chhorā chup hai.} — The teacher makes a sign and the boys are quiet.
  \item \emph{Recall:} say the Marwadi for agreed, then the Marwadi for silent, then say \emph{isāro} again
\end{itemize}

\begin{tcolorbox}[breakable,colback=teal!4,colframe=teal!35!black,title={Before you move on}]
What does \mw{इसारो} mean? (\textquotedblleft{}A gesture\textquotedblright{}.) Say it once more.
\end{tcolorbox}
\section[\texorpdfstring{\mw{भेद} (bhed)}{भेद (bhed)}]{\mw{भेद} (bhed) — a secret}

\label{lesson:MW-C273-bhed}



\begin{quote}
Before the new one: say the Marwadi for silent, then the Marwadi for a gesture.
\end{quote}

\subsection*{You'll want to know: \mw{भेद}}
\textbf{\mw{भेद}} — \emph{bhed} — \textquotedblleft{}a secret\textquotedblright{}.

\begin{quote}
\mw{यो} \mw{भेद} \mw{है}, \mw{चुप} \mw{रैवो।}
\end{quote}

\begin{quote}
\emph{yo bhed hai, chup raivo.}
\end{quote}

\begin{quote}
This is a secret: keep quiet.
\end{quote}

A masculine noun.

\begin{cousinweb}
\textbf{\mw{भेद}} (\emph{bhed}) is a tatsama, taken from Sanskrit \emph{bheda}, \textquotedblleft{}breaking open\textquotedblright{}, and so \textquotedblleft{}disclosing (a secret)\textquotedblright{} (Monier-Williams). Here it names the secret itself, a sense Monier-Williams does not list.
\end{cousinweb}

\begin{grammarlens}[title={What you've built}]
One more word for asking and answering.
\end{grammarlens}

\subsection*{Your turn}
\begin{itemize}
\raggedright
  \item \emph{Say it:} \emph{bhed}
  \item \emph{Say it:} \emph{bhed}, once more
  \item \emph{Say it:} \emph{yo bhed hai, chup raivo.} — This is a secret: keep quiet.
  \item \emph{Recall:} say the Marwadi for silent, then the Marwadi for a gesture, then say \emph{bhed} again
\end{itemize}

\begin{tcolorbox}[breakable,colback=teal!4,colframe=teal!35!black,title={Before you move on}]
What does \mw{भेद} mean? (\textquotedblleft{}A secret\textquotedblright{}.) Say it once more.
\end{tcolorbox}
\section[\texorpdfstring{\mw{छिपावणो} (chhipāvṇo)}{छिपावणो (chhipāvṇo)}]{\mw{छिपावणो} (chhipāvṇo) — to hide (something); to keep secret}

\label{lesson:MW-C273-chhipavno}



\begin{quote}
Before the new one: say the Marwadi for a gesture, then the Marwadi for a secret.
\end{quote}

\subsection*{You'll want to know: \mw{छिपावणो}}
\textbf{\mw{छिपावणो}} — \emph{chhipāvṇo} — \textquotedblleft{}to hide (something); to keep secret\textquotedblright{}.

\begin{quote}
\mw{वो} \mw{भेद} \mw{छिपावै} \mw{है।}
\end{quote}

\begin{quote}
\emph{vo bhed chhipāvai hai.}
\end{quote}

\begin{quote}
He keeps the secret.
\end{quote}

Past (I …): \mw{म्हैं} \mw{भेद} \mw{छिपायो}, \emph{mhaĩ bhed chhipāyo}.

\begin{cousinweb}
\textbf{\mw{छिपावणो}} (\emph{chhipāvṇo}) is inherited from a reconstructed \emph{*chipp}, \textquotedblleft{}cover, hide\textquotedblright{}; CDIAL lists Old Marwari \emph{chipāvaï}, \textquotedblleft{}hides\textquotedblright{} (CDIAL 4994).
\end{cousinweb}

\begin{grammarlens}[title={What you've built}]
One more word for asking and answering.
\end{grammarlens}

\subsection*{Your turn}
\begin{itemize}
\raggedright
  \item \emph{Say it:} \emph{chhipāvṇo}
  \item \emph{Say it:} \emph{chhipāvṇo}, once more
  \item \emph{Say it:} \emph{vo bhed chhipāvai hai.} — He keeps the secret.
  \item \emph{Recall:} say the Marwadi for a gesture, then the Marwadi for a secret, then say \emph{chhipāvṇo} again
\end{itemize}

\begin{tcolorbox}[breakable,colback=teal!4,colframe=teal!35!black,title={Before you move on}]
What does \mw{छिपावणो} mean? (\textquotedblleft{}To hide (something); to keep secret\textquotedblright{}.) Say it once more.
\end{tcolorbox}
\section[\texorpdfstring{\mw{वचन} (vachan)}{वचन (vachan)}]{\mw{वचन} (vachan) — a promise}

\label{lesson:MW-C273-vachan}



\begin{quote}
Before the new one: say the Marwadi for a secret, then the Marwadi for to keep secret.
\end{quote}

\subsection*{You'll want to know: \mw{वचन}}
\textbf{\mw{वचन}} — \emph{vachan} — \textquotedblleft{}a promise\textquotedblright{}.

\begin{quote}
\mw{यो} \mw{म्हारो} \mw{वचन} \mw{है।}
\end{quote}

\begin{quote}
\emph{yo mhāro vachan hai.}
\end{quote}

\begin{quote}
This is my promise.
\end{quote}

A masculine noun.

\begin{cousinweb}
\textbf{\mw{वचन}} (\emph{vachan}) is a tatsama, taken from Sanskrit \emph{vacana}, \textquotedblleft{}speaking; a word\textquotedblright{} (Monier-Williams), which does not list the sense \textquotedblleft{}a promise\textquotedblright{}. CDIAL gives no Marwari form under \emph{vacana}, but lists the inherited Gujarati \emph{veṇ}, \textquotedblleft{}word, promise\textquotedblright{} (CDIAL 11199): a word given is a promise.
\end{cousinweb}

\begin{grammarlens}[title={What you've built}]
One more word for asking and answering.
\end{grammarlens}

\subsection*{Your turn}
\begin{itemize}
\raggedright
  \item \emph{Say it:} \emph{vachan}
  \item \emph{Say it:} \emph{vachan}, once more
  \item \emph{Say it:} \emph{yo mhāro vachan hai.} — This is my promise.
  \item \emph{Recall:} say the Marwadi for a secret, then the Marwadi for to keep secret, then say \emph{vachan} again
\end{itemize}

\begin{tcolorbox}[breakable,colback=teal!4,colframe=teal!35!black,title={Before you move on}]
What does \mw{वचन} mean? (\textquotedblleft{}A promise\textquotedblright{}.) Say it once more.
\end{tcolorbox}
